\section*{Chapter \ref{ch:dv:implied-volatility-and-its-surface} --- Implied Volatility and Its Surface}

\begin{solution}{exo:dv:implied-volatility-and-its-surface:1}
$k=\ln(90/100.37)=-0.1091$; $w=0.2047^2\times91/365=0.0104$.
\end{solution}

\begin{solution}{exo:dv:implied-volatility-and-its-surface:2}
$20.47-14.19=6.27$ volatility points. At one year the same strikes are closer to the money in
standard deviations ($\ln(K/F)/\sqrt T$ is halved), and the skew per unit of $k$ is itself
smaller, decaying roughly like $T^{-1/2}$: the one-year 90--110 skew is 3.3 points.
\end{solution}

\begin{solution}{exo:dv:implied-volatility-and-its-surface:3}
$w$ is $0.2^2\times0.25=0.0100$ at three months and $0.13^2\times0.5=0.0085$ at six: total
variance falls, a \omterm{def:dv:implied-volatility-and-its-surface:static}{calendar arbitrage}. Sell the three-month straddle (or call) and buy the
six-month one at the same forward-moneyness: the long-dated option is worth at least the
short-dated one at the first expiry.
\end{solution}

\begin{solution}{exo:dv:implied-volatility-and-its-surface:4}
Without the skew: $P\Phi(d_2)=0.4980$. The correction is $-\mathcal V\,\partial\sigma/\partial K
=19.77\times0.0031=0.0604$ (per unit of volatility): 0.5584 with the skew.
\end{solution}

\begin{solution}{exo:dv:implied-volatility-and-its-surface:5}
\omterm{def:dv:implied-volatility-and-its-surface:sticky}{Sticky strike}: the Black--Scholes delta, 0.533. \omterm{def:dv:implied-volatility-and-its-surface:sticky}{Sticky delta}: $0.533+19.77\times0.00306=0.593$.
Six shares per 100 calls.
\end{solution}

\begin{solution}{exo:dv:implied-volatility-and-its-surface:6}
An index's risk-neutral distribution adds the members' tendency to fall together: in a crash
correlations rise, so the index's left tail is fatter than any member's relative to its own
volatility. Individual risk-neutral distributions are far less negatively skewed than the
index's; the gap is the correlation skew of chapter 17.
\end{solution}

\begin{solution}{exo:dv:implied-volatility-and-its-surface:7}
Six butterfly violations, the worst at $g=-2.06$ at $k=-0.1$. A downward error of 6 points on
the six-month at-the-money quote leaves $w$ increasing; 6.5 points (the first half-point step
that does) creates one calendar violation, 7 points one, 8 points three.
\end{solution}

\begin{solution}{exo:dv:implied-volatility-and-its-surface:8}
Implied volatility at a strike is a price, not the volatility expected if the index reaches that
strike. The 80 put's 27 prices the probability and severity of a fall in one number, including
a premium for insurance and for jumps, which no diffusion volatility at 80 describes. What the
market would expect conditional on a fall is the local volatility of chapter 9 or the model's
dynamics, not the implied volatility at that strike.
\end{solution}

\begin{solution}{pb:dv:implied-volatility-and-its-surface:1}
\textbf{1.} Forward 100.37; discount factor 0.99255.
\textbf{2.} 16.36\% and 26.78\%.
\textbf{3.} $\ln(80/100.37)=-0.2269$.
\textbf{4.} 0.219.
\textbf{5.} 0.006.
\textbf{6.} 0.31\%.
\textbf{7.} 1.89\%.
\textbf{8.} The probability is minus the strike derivative of the put divided by the discount
factor, and the put's volatility changes with the strike: the skew adds $\mathcal V\,
\partial\sigma/\partial K$, which $\Phi(-d_2)$ at a fixed volatility ignores.
\textbf{9.} It is the forward value of a digital put struck at 80, and the limit of a put spread
$(P(80+h)-P(80-h))/(2h)$ divided by the discount factor.
\textbf{10.} From the centre and the right tail into the left tail (and a little into the region
just above the forward, where the peak rises).
\textbf{11.} No: it is a price, the probability under the measure that makes prices expectations,
and includes the premium investors pay to avoid the fall.
\textbf{12.} The risk premium for crash insurance, the demand of hedgers, and the dealers' cost of
carrying the risk.
\textbf{13.} The put gains from the fall in both regimes; under \omterm{def:dv:implied-volatility-and-its-surface:sticky}{sticky strike} its volatility is
unchanged at 26.8, under \omterm{def:dv:implied-volatility-and-its-surface:sticky}{sticky delta} it falls, since the put is now closer to the money
(23.8 at the new moneyness): \omterm{def:dv:implied-volatility-and-its-surface:sticky}{sticky strike} gives the larger gain.
\textbf{14.} The crash showed that a fall of a fifth within two days was possible; buyers of
protection and sellers who had lost repriced the left tail, and have kept doing so.
\textbf{15.} A fall is sudden: a jump or a burst of volatility is a larger share of the variance
to a short expiry than to a long one, so short-dated skews are steeper (chapters 12--13).
\textbf{16.} That the market charges about six times the flat-volatility probability, and that
the price is 0.22 per 100 of notional per quarter, about 0.9\% a year if rolled.
\textbf{17.} A put spread (buy 80, sell 70) or a collar: cheaper, but the protection stops below
70, or the upside is given away.
\textbf{18.} The skewed surface, which is where the protection can actually be bought; the flat
surface understates the tail by a factor of six.
\textbf{19.} 0.31\% flat, 1.89\% skewed: six times more.
\textbf{20.} The market's price of the distribution's left tail, above what one volatility can
describe.
\end{solution}

\begin{solution}{iq:dv:implied-volatility-and-its-surface:1}
Investors buy index puts as insurance, crashes are downward and sudden, and volatility rises
when prices fall; before 1987 the market priced near-lognormal tails, and the crash showed
that a fall of a fifth in a day or two was possible. Sellers of puts have charged for it since.
\iqlookfor{demand for insurance, jumps, the leverage effect, and the historical break.}
\end{solution}

\begin{solution}{iq:dv:implied-volatility-and-its-surface:2}
Breeden--Litzenberger: $q(K)=\partial^2C/\partial K^2/P(0,T)$. In practice fit an
arbitrage-free smile first (raw quotes are too noisy to differentiate twice), then differentiate
the smooth call prices; check that the density integrates to one and reprices the forward.
\iqlookfor{the formula and the need to smooth first.}
\end{solution}

\begin{solution}{iq:dv:implied-volatility-and-its-surface:3}
\omterm{def:dv:implied-volatility-and-its-surface:sticky}{Sticky strike}: each strike keeps its volatility, so the new at-the-money volatility is read
higher up the skew (18.1 from 16.4 in the chapter) and delta is the Black--Scholes delta.
\omterm{def:dv:implied-volatility-and-its-surface:sticky}{Sticky delta}: the smile moves with the forward, at-the-money volatility is unchanged, and
delta adds \omterm{def:dv:greeks-and-the-hedging-pnl:greeks}{vega} times the slope, higher for a negative skew.
\iqlookfor{both the level change and the hedge change.}
\end{solution}

\begin{solution}{iq:dv:implied-volatility-and-its-surface:4}
Calls decreasing in strike and with slope above $-P(0,T)$; convex in strike (butterflies
non-negative), equivalently $g(k)\ge0$ in total variance; total variance non-decreasing in
expiry at fixed forward-moneyness. Check each slice on a fine grid of $k$ and each pair of
consecutive slices; with bid and ask, check at executable prices.
\iqlookfor{the three conditions and their total-variance form.}
\end{solution}

\begin{solution}{iq:dv:implied-volatility-and-its-surface:5}
It uses a lognormal tail, ignoring the skew. The probability is the strike derivative of the
put, which includes the skew term; on a typical index surface the error at 20\% below is a
factor of five or more (six in the chapter).
\iqlookfor{the skew in the derivative, and the size of the error.}
\end{solution}

\begin{solution}{iq:dv:implied-volatility-and-its-surface:6}
Work in total variance and \omterm{def:dv:implied-volatility-and-its-surface:surface}{log-moneyness}; within an expiry use a smooth parametrisation or a
spline in $k$ that respects $g\ge0$ (SVI, chapter 8), with care at the ends; between expiries
interpolate $w$ linearly in time at fixed $k$, which preserves calendar monotonicity; handle
dividends and events in the time coordinate; check the result.
\iqlookfor{total variance, monotone interpolation in time, and a check.}
\end{solution}
